\documentclass[11pt]{article}
\usepackage[a4paper,margin=25mm]{geometry}
\usepackage{amsmath,amssymb,amsthm,booktabs,array,longtable}
\usepackage[T1]{fontenc}
\usepackage{lmodern}
\usepackage[hidelinks]{hyperref}
\emergencystretch=2em
\newtheorem{theorem}{Theorem}
\newtheorem{proposition}[theorem]{Proposition}
\newtheorem{corollary}[theorem]{Corollary}
\newcommand{\Z}{\mathbb Z}
\newcommand{\Q}{\mathbb Q}
\newcommand{\F}{\mathbb F}
\title{Sharp tails for Bala's inverse-function conjectures\\
and a specialization of Su's period theorem}
\author{Haobo Ma\\\small ChronoAI Pte Ltd, Singapore
\and Wenlin Zhang\\\small National University of Singapore\\[-1pt]
\small ChronoAI Pte Ltd, Singapore}
\date{November 14, 2026}
\hypersetup{pdftitle={Sharp tails for Bala's inverse-function conjectures and a specialization of Su's period theorem},pdfauthor={Haobo Ma and Wenlin Zhang}}
\begin{document}
\maketitle
\begin{abstract}
For Peter Bala's totient-square conjecture in OEIS A281184, comment of
February 16, 2026, revision 41, we specialize Lien-Hung Su's Theorem 3.2
and Corollary 3.3 (October 31, 2026) to obtain the asserted eventual
period by identifying the coefficients with integer derivation iterates.
For Bala's related examples in A281428, March 3, 2026, revision 54,
we give finite-dimensional certificates proving the least modulo-$27$
tail period $18$ from the exact index $2$, and divisibility by $385$
from the exact index $6$.
The former improves Su's sufficient index $3$.
We also reprove seven companion least periods, five row-multiplier
identities, and two coefficient identities, distinguishing Su's general
theorems, written proofs, exact finite certificates, and the scope of
the linked Lean declarations. No priority claim is made.
\end{abstract}

\section{The coefficients and the sharp tails}
Bala's question in \cite{oeis184} concerns the even exponential
coefficients of $C^3$, where $S'=C^4$, $C'=SC^3$, $S(0)=0$, $C(0)=1$.
Su \cite{su} proves the general modular assertions for this flow and its
neighbors. The purpose here is to explain that identification and give
explicit small-rank certificates, with sharp tail origins for the
odd-flow examples. Su's general prime-power argument supplies all-modulus
conclusions; the finite-dimensional argument below supplies exact periods
and origins for specified moduli at the cost of rank $\lceil m/2\rceil$.
The differential descriptions themselves are inherited from the OEIS
contributors and Bala's account of inverse functions \cite{graves}.

Fix an integer $b\geq2$. In
\[
 R=\Z[s,c]/(c^2-s^2-1),\qquad
 D_b s=c^b,\quad D_b c=sc^{b-1},\quad E(f)=f(0,1),
\]
the derivation is well defined because it kills $c^2-s^2-1$.
Define
\begin{equation}\label{eq:coeff}
 A_b(n,e)=E(D_b^{2n}c^e)=(2n)![x^{2n}]C_b(x)^e
 \quad(n,e\geq0),
\end{equation}
where $S_b'=C_b^b$, $C_b'=S_bC_b^{b-1}$ and $(S_b(0),C_b(0))=(0,1)$.
To justify the equality, the iterated Leibniz rule shows that
$f\mapsto\sum_{j\geq0}E(D_b^jf)x^j/j!$ is multiplicative and intertwines
$D_b$ with differentiation. Its images of $s,c$ therefore solve this
system; coefficient recursion gives uniqueness in $\Q[[x]]$.
It also proves integrality of every $A_b(n,e)$.
Uniqueness under $x\mapsto-x$ proves that $S_b$ is odd and $C_b$ even.
We have $C_b^2-S_b^2=1$ and
$(\log(C_b+S_b))'=C_b^{b-1}$; hence the $b=4,e=3$ definition is
exactly A281184. Put $u(n)=A_5(n,1)$, which is A281428.
Here and below a least tail period is the least positive period of the
eventually periodic residue sequence, and an exact origin is its earliest
valid starting index in the source's indexing.

\begin{theorem}[Sharp origins for Bala's odd-flow examples]\label{thm:sharp}
The sequence $u(n)$ has least eventual period $18$ modulo $27$,
with exact origin $n=2$. Moreover
\[
 u(n)\equiv0\pmod{385}\quad(n\geq6),\qquad u(5)\equiv105\pmod{385}.
\]
Its exact zero-tail origins modulo $5,7,11$ are $3,4,6$, respectively.
\end{theorem}
The all-modulus statements underlying these examples are already proved
by Su, Theorem 4.1. In particular his $3$-primary bound gives origin $3$
modulo $27$; Theorem~\ref{thm:sharp} sharpens it to $2$.
Bala's displayed sufficient origin $7$ is valid and is not an assertion
of minimality \cite[March 3, 2026 example block]{oeis428}.

\section{Finite certificates and their all-index meaning}
Write $z=s^2$, $\eta\equiv e\pmod2$, $\eta\in\{0,1\}$, and
$f_0=(1+z)^{(e-\eta)/2}$. Direct differentiation gives
\begin{equation}\label{eq:operator}
 D_b^2(c^\eta f)=c^\eta B_{b,\eta}f,\qquad
 \begin{aligned}
 g&=\eta f+2(1+z)f',\\
 B_{b,\eta}f&=(1+z)^{b-2}
       \bigl((1+(\eta+b-1)z)g+2z(1+z)g'\bigr).
 \end{aligned}
\end{equation}
Thus $A_b(n,e)=[z^0]B_{b,\eta}^{\,n}f_0$ over the integers.
Fix $m\geq2$ and $h=\lceil m/2\rceil$. The coefficient of $z^{j-1}$
in $B_{b,\eta}(z^j)$ is $2j(2j-1)$ and all other terms have degree
at least $j$. At $j=h$ that coefficient is divisible by $m$, since
$m=2h$ or $2h-1$. For $j>h$ even the degree-lowering term has degree
at least $h$. Therefore $B_{b,\eta}$ induces an endomorphism $M$ of the
free rank-$h$ module $(\Z/m\Z)[z]/(z^h)$, and
\begin{equation}\label{eq:scalar}
 A_b(n,e)\equiv\ell M^nf_0\pmod m,\qquad \ell(f)=[z^0]f.
\end{equation}

\begin{proposition}[Propagation of a finite scalar certificate]\label{prop:ch}
For $P\geq1$, $r\in\Z/m\Z$, and $t\geq0$, if
$\ell M^j(M^P-rI)f_0=0$ for $t\leq j<t+h$, then
$A_b(j+P,e)=rA_b(j,e)$ modulo $m$ for every $j\geq t$.
Likewise $h$ consecutive zero values beginning at $t$ imply a zero tail.
\end{proposition}
\begin{proof}
The characteristic polynomial of $M$ is monic of degree $h$.
Cayley--Hamilton over the commutative ring $\Z/m\Z$ gives the same
order-$h$ scalar recurrence for $\ell M^jv$ for every state $v$.
Induction propagates $h$ consecutive zeros. Take
$v=(M^P-rI)f_0$ for the first claim and $v=f_0$ for the second.
\end{proof}

\begin{proof}[Proof of Theorem~\ref{thm:sharp}]
Take $(b,e,m,h)=(5,1,27,14)$, so $f_0=1$.
Formula~\eqref{eq:operator}, with all coefficients reduced modulo $27$,
gives
\[
 (u(j+18)-u(j))_{j=2}^{15}=(0,0,0,0,0,0,0,0,0,0,0,0,0,0).
\]
These are exactly fourteen premises of Proposition~\ref{prop:ch}.
For transparency the first period, at indices $2$ through $19$, is
\[
 (17,1,8,19,17,1,26,19,17,10,26,19,8,10,26,1,8,10).
\]
The pairs $u(2)=17$, $u(11)=10$ and $u(8)=26$ exclude periods
$9$ and $6$. The least eventual period divides $18$: translate any
eventual period back into this $18$-periodic tail and use the common-divisor
property for periods of a periodic word. Every proper divisor of $18$
divides $9$ or $6$, so the least period is $18$.
Also $u(1)=1\ne10=u(19)$ modulo $27$, proving the exact origin.

For the zero tails, the same operator gives the following full premises;
the column headed ``zeros'' lists $h$ consecutive values from $t$.
\[
\begin{array}{@{}rrrrl@{}}\toprule
 m&h&t&u(t-1)\bmod m&\text{zeros}\\\midrule
 5&3&3&2&(0,0,0)\\
 7&4&4&4&(0,0,0,0)\\
 11&6&6&6&(0,0,0,0,0,0)\\\bottomrule
\end{array}
\]
Proposition~\ref{prop:ch} proves the zero tails; the nonzero preceding
values prove their exact origins. Chinese remaindering gives the $385$
congruence from $6$. Directly the residues at indices $0$ through $6$
modulo $385$ are $(1,1,17,95,315,105,0)$, proving sharpness.
All displayed finite premises use only the explicit polynomial rule
\eqref{eq:operator}; \texttt{verify.py} supplies their exact arithmetic.
\end{proof}

\section{The companion examples and two-dimensional rows}
The source identifications, including the one-index offset, are
\begin{equation}\label{eq:family}
\begin{aligned}
 \mathrm{A281180}(n)&=A_4(n-1,4) &&(n\geq1),\\
 \mathrm{A281181}(n)&=A_4(n,1),&
 \mathrm{A281183}(n)&=A_4(n,2),\\
 \mathrm{A281184}(n)&=A_4(n,3),&
 \mathrm{A281440}(n)&=A_8(n,1).
\end{aligned}
\end{equation}
The first follows by differentiating $S_4$, the others from their e.g.f.s.
Su's Proposition 3.6 proves the next seven least periods. We give a
common certificate so the row observations have a precise interpretation.

\begin{corollary}[Companion periods, reproved]\label{cor:periods}
The table gives least periods $P$ and exact origins $n_0$ in source
indices. In each row $h$ scalar period differences from $n_0$ are zero.
\[
\begin{array}{@{}lrrrr@{}}\toprule
 \text{sequence}&m&P&n_0&h\\\midrule
 \mathrm{A281180}&7&18&1&4\\
 \mathrm{A281181}&10&8&0&5\\
 \mathrm{A281181}&9&3&1&5\\
 \mathrm{A281181}&13&72&0&7\\
 \mathrm{A281183}&11&25&1&6\\
 \mathrm{A281184}&11&25&1&6\\
 \mathrm{A281440}&13&72&0&7\\\bottomrule
\end{array}
\]
Also $\mathrm{A281181}(n)\equiv1\pmod{12}$ for every $n\geq0$.
\end{corollary}
\begin{proof}
Use \eqref{eq:family} in \eqref{eq:operator} and Proposition~\ref{prop:ch}.
Here are the witnesses excluding the maximal proper divisor shifts;
each entry is $(n:a(n),a(n+P/q))$ for the prime divisor $q$ of $P$.
\[
\begin{array}{@{}lll@{}}\toprule
 \text{sequence, modulus}&q&\text{witness}\\\midrule
 \mathrm{A281180},7&2,3&(1:1,6),\ (1:1,2)\\
 \mathrm{A281181},10&2&(0:1,9)\\
 \mathrm{A281181},9&3&(1:1,4)\\
 \mathrm{A281181},13&2,3&(0:1,12),\ (0:1,9)\\
 \mathrm{A281183},11&5&(1:2,10)\\
 \mathrm{A281184},11&5&(1:3,4)\\
 \mathrm{A281440},13&2,3&(0:1,12),\ (0:1,9)\\\bottomrule
\end{array}
\]
For rows with $n_0=1$ and source domain starting at $0$, the preceding
pairs $(a(0),a(P))$ are respectively $(1,7),(1,10),(1,3)$ for
A281181 modulo $9$, A281183 and A281184 modulo $11$.
The other origins are the beginning of the source domain.
For the last claim use $(b,e,m,h)=(4,1,12,6)$: the seven values at
$0\leq n\leq6$ are all $1$ modulo $12$. Their six successive
differences and Proposition~\ref{prop:ch} prove the constant residue.
\end{proof}

\begin{corollary}[The five row multipliers, reproved]\label{cor:rows}
The following identities hold for all $n\geq n_0$:
$a(n+r)\equiv\lambda a(n)\pmod m$.
\[
\begin{array}{@{}lrrrr@{}}\toprule
 \text{sequence}&m&r&\lambda&n_0\\\midrule
 \mathrm{A281180}&7&6&2&1\\
 \mathrm{A281181}&13&6&7&0\\
 \mathrm{A281183}&11&5&5&1\\
 \mathrm{A281184}&11&5&5&1\\
 \mathrm{A281440}&13&6&6&0\\\bottomrule
\end{array}
\]
\end{corollary}
\begin{proof}
For each row \eqref{eq:operator} gives $h=\lceil m/2\rceil$ zero
differences $a(n+r)-\lambda a(n)$ from $n_0$. Apply
Proposition~\ref{prop:ch}, with the offset in \eqref{eq:family}.
Thus the rows in the source rectangles form geometric progressions
columnwise. Su's Proposition 3.4 gives a general higher-order row
relation for $b=4$, reducing to a unit multiplier for squarefree moduli
after its stated threshold. It accounts for the four $b=4$ rows;
the certificates here also establish the listed origins.
The $b=8$ row follows directly from its certificate, and is consistent
with Su's prime-characteristic operator law in Section 2.
\end{proof}

\section{The all-modulus theorems and coefficient siblings}
For an even flow exponent $b\geq2$, Su's Theorem 3.2 states, for $m\geq3$,
\begin{equation}\label{eq:su}
 A_b(n+\varphi(m)^2/2,e)\equiv A_b(n,e)\pmod m
 \quad(n\geq N_b(m),\ e\geq0).
\end{equation}
Here $N_b(m)=\max_{p^r\Vert m}N_b(p^r)$, where $N_b(2^r)=r$ and,
for odd $p$,
\[
 \alpha_{b,p}=(-1)^{(p-1)/2}
 \prod_{j=0}^{(p-3)/2}\frac{b+2j}{2(j+1)}\ \text{in }\F_p,\qquad
 N_b(p^r)=\begin{cases}
 r,&\alpha_{b,p}\ne0,\\
 \lceil r(1+(p-1)p^{r-1})/2\rceil,&\alpha_{b,p}=0.
 \end{cases}
\]
The denominators are units. For $b=4$, the even numerator factors
$4,6,\ldots,p+1$ contain no multiple of odd $p$, so
$N_4(m)=\max_{p\mid m}v_p(m)$. Vanishing branches can occur for $b=8$;
the stated bound retains them. Applying \eqref{eq:su} twice proves
Bala's five totient-square assertions via \eqref{eq:family}.
For A281180 the sufficient source index is $N_4(m)+1$.
At $m=1$ all congruences hold, and at $m=2$ each column is constant
from $1$: the elementary recurrence
\begin{equation}\label{eq:rec}
 A_b(n,e)=e(e+b-1)A_b(n-1,e+2b-2)
          -e(e+b-2)A_b(n-1,e+2b-4),\quad A_b(0,e)=1
\end{equation}
shows that the tail is $0$ for even $e$ and $1$ for odd $e$.
This follows by differentiating $c^e$ twice, and for $b=4$ is
Manyama's source recurrence \cite{oeis184}.

Su's Theorem 4.1 proves both conjectures in A281428, revision 54,
Bala's March 3, 2026 comment block:
\begin{itemize}
\item If $(m,6)=1$, then $u(n)$ is eventually zero modulo $m$.
For $p\geq5$, $r\geq1$, a sufficient index modulo $p^r$ is
$\lceil r(1+(p-1)p^{r-1})/2\rceil$.
\item For every $m\geq1$, $\varphi(m)$ is an eventual period of $u$
modulo $m$. The source poses this for moduli divisible by $2$ or $3$;
the extension to all $m$ follows by joining the zero branches.
More precisely $u(n+2^t)\equiv u(n)\pmod{2^{t+2}}$ for $n,t\geq0$,
and $u(n+2\cdot3^{r-1})\equiv u(n)\pmod{3^r}$ for $n\geq r\geq1$.
\end{itemize}
These are cited general theorems, with the coefficient identification
\eqref{eq:coeff}; Theorem~\ref{thm:sharp} provides the sharper finite-modulus
origins by a separate written certificate argument.

\begin{corollary}[Full-index companion and derivative identities]\label{cor:identities}
For A281182 set $t(j)=j![x^j](C_4+S_4)$. Then
$t(2n)=A_4(n,1)$ and $t(2n+1)=A_4(n,4)$.
Consequently $\varphi(m)^2$ is a full-index tail period for $m\geq3$
from index $2N_4(m)$, and $2\varphi(m)^2$ is a tail period for all $m\geq1$.
The full-index shift $1$ fails modulo $2$; it is not asserted in this source.

For A281183, let $v(y)=1+y^2$ and
$T(f)=v\,d(vf)/dy$. Then $(T^{2n}1)(0)=A_4(n,2)$.
For A281181, $(S_4/C_4)'=C_4$, and the $x^5$ coefficient of $S_4/C_4$
is $13/5!$.
\end{corollary}
\begin{proof}
Parity and $S_4'=C_4^4$ give the two coefficient equalities. Doubling
the even-column shift in \eqref{eq:su} preserves each parity class.
Modulo $2$ the even terms are $1$, while the odd terms are $0$ from
$j=3$, so a shift of $2$ works but $1$ does not.
For the second assertion set $L=v^2d/dy$. Then $T=v^{-1}Lv$ and
$T^j1=v^{-1}L^jv$ in $\Q[[y]]$. Evaluation at $y=0$, using the flow
$S_4'=v(S_4)^2$, gives $(L^jv)(0)=j![x^j]C_4^2$.
This proves Bala's tentative formula of December 22, 2025;
the same proof is Su's Proposition 3.5.
Finally $(S_4/C_4)'=C_4(C_4^2-S_4^2)=C_4$.
Since $A_4(2,1)=13$, integration gives $13x^5/5!$.
Thus the printed denominator $6!$ in the December 14, 2025 expansion
in A281181 should be $5!$; the derivative identity itself is correct.
\end{proof}

\begin{corollary}[Weighted-derivative normalization]\label{cor:factorial}
In A281428's January 22, 2026 formula, let
$f_k(0,\theta)=\sec^{k-3}\theta$ and
$f_k(j,\theta)=\partial_\theta(\sec^3\theta f_k(j-1,\theta))$.
For integer $k\geq0$ and $j\geq0$, the correctly normalized identity is
$f_k(j,0)=j![x^j]C_5(x)^k$.
\end{corollary}
\begin{proof}
Put $L=\sec^3\theta\,\partial_\theta$, $w=\sec^3\theta$.
The recurrence operator is $w^{-1}Lw$, so
$f_k(j,\theta)=w^{-1}L^j(\sec^k\theta)$.
Under $\theta'=\sec^3\theta$, $\theta(0)=0$, the series
$\sec\theta(x)=C_5(x)$ and $\tan\theta(x)=S_5(x)$ solve the $b=5$ flow.
Evaluation therefore gives the asserted derivative coefficient.
At $k=1,j=2$ the left side is $1$, whereas $[x^2]C_5=1/2$.
Su's remark after Theorem 4.1 states the same correction.
\end{proof}

\section{Coverage and verification scope}
The following table separates source assertions, examples, and extensions.
The OEIS revision numbers identify the versions consulted; dates locate
the signed source blocks, not claims of priority.
\small
\begin{longtable}{@{}>{\raggedright\arraybackslash}p{.43\textwidth}>{\raggedright\arraybackslash}p{.52\textwidth}@{}}
\toprule Source item and exact locator & Coverage\\\midrule\endhead
A281184 rev.~41, Bala comment, Feb.~16, 2026;
A281180 rev.~86, Feb.~18, 2026;
A281181 rev.~103 and A281183 rev.~60, Feb.~16, 2026;
A281440 rev.~17, Feb.~19, 2026 &
Proved: coefficient identifications and specializations of Su's Theorem 3.2
and Corollary 3.3. Prior proof cited: Su (2026).\\
Seven displayed least-period examples in these entries; A281181 Curtz
comments, Jan.~20, 2017 &
Reproved; prior proof Su (2026), Proposition 3.6.
Corollary~\ref{cor:periods}; the mod-$12$ difference assertion is also proved.\\
Five rectangular displays and Bala's ``two-dimensional'' remarks in
the dated modular example blocks &
Proved as the five precise multiplier identities of
Corollary~\ref{cor:rows}; prior general row relation for $b=4$:
Su (2026), Proposition 3.4; the $b=8$ certificate is stated separately.
Unspecified stronger meanings are not covered.\\
A281428 rev.~54, Bala conjectures 1 and 2, Mar.~3, 2026 &
Proved by the cited Su (2026), Theorem 4.1, with all-modulus scope stated above.\\
A281428, same-date examples, moduli $27$ and $385$ &
Proved, with exact origins $2$ and $6$ in Theorem~\ref{thm:sharp}.
Su's general mod-$27$ origin is sufficient but not sharp.\\
A281183 rev.~60, tentative $g$ formula, Dec.~22, 2025 &
Reproved; prior proof Su (2026), Proposition 3.5;
Corollary~\ref{cor:identities}.\\
A281181 rev.~103, $S/C$ expansion, Dec.~14, 2025 &
Refuted-and-corrected coefficient $13/6!$ to $13/5!$;
the asserted derivative identity is proved.\\
A281428 rev.~54, generalized derivative formula, Jan.~22, 2026 &
Refuted-and-corrected ordinary-coefficient normalization;
Corollary~\ref{cor:factorial}. Prior correction cited: Su (2026), remark after Theorem 4.1.\\
A281182 rev.~32, defining full-index e.g.f. &
Not posed; full-index periods proved in Corollary~\ref{cor:identities}.
The universal shift-$1$ reading modulo $2$ is not asserted in the source;
it fails on the alternating tail.\\
Other cross-referenced sequences and unrelated inverse-function formulas &
Not covered; no claim of resolving their separate problems.\\\bottomrule
\end{longtable}
\normalsize

\begin{samepage}
The following pinned Lean declarations state the formal scope.
\begin{itemize}
\setlength{\itemsep}{0pt}
\setlength{\parsep}{0pt}
\item
\href{https://github.com/the-omega-institute/trureturing/blob/92de86888730265f84e24d28bdb67546c543a966/D5/S1/Recurrence/Periodic/BalaA281184TotientSquarePeriod.lean}{\texttt{BalaA281184TotientSquarePeriod.lean}}
\par\noindent
The declaration \texttt{result} in its matching namespace states existence of the
source series, the integer even coefficients of $C^3$, and eventual
shift $\varphi(m)^2$ for every positive $m$.
\end{itemize}
\end{samepage}
\begin{itemize}
\setlength{\itemsep}{0pt}
\setlength{\parsep}{0pt}
\item
\href{https://github.com/the-omega-institute/trureturing/blob/92de86888730265f84e24d28bdb67546c543a966/D5/S1/Recurrence/Periodic/BalaA281181TotientSquarePeriod.lean}{\texttt{BalaA281181TotientSquarePeriod.lean}},
\par\noindent
The namespace \texttt{BalaPolynomialSystem} has
\texttt{even\_sequence\_coefficient}, identifying the normalized coefficients,
and \texttt{even\_sequence\_period}, giving the square-period shift for
$q=3$ or $7$, arbitrary observable exponent, and flow exponent $b=q+1$.
\item
\href{https://github.com/the-omega-institute/trureturing/blob/92de86888730265f84e24d28bdb67546c543a966/D5/S1/Recurrence/Residue/BalaQuarticEGFEventualDivisibility.lean}{\texttt{BalaQuarticEGFEventualDivisibility.lean}}
\par\noindent
The declaration \texttt{result} in its matching namespace states existence of a
normalized source solution and eventual integer-multiple divisibility
of its even exponential coefficients for every positive modulus coprime
to $6$. It states neither the exact $385$ origin nor the second
A281428 conjecture about periods at moduli divisible by $2$ or $3$.
\end{itemize}
All three files are pinned to the same revision in their links.
The half-square shifts and general odd-flow conclusions are the cited
Su theorems. Exact origins, the eight least-period certificates, five
row relations, and normalization identities have the written proofs above.

The supplementary \texttt{verify.py} and \texttt{checks.json} contain the
eight complete finite period certificates, their minimality witnesses,
five row certificates, three zero-tail certificates, and the mod-$12$
certificate. They also compare the integer recurrence \eqref{eq:rec}
and polynomial operator \eqref{eq:operator} at all $257$ indices
$0\leq n\leq256$ for A281184 and at $65$ indices for eight further
parameter pairs. Fourteen additional moduli have $64$ finite comparisons
each. Those coefficient and modular comparisons are finite checks;
the all-index consequences use Proposition~\ref{prop:ch} or the cited
all-modulus theorems. All arithmetic is exact, with no rounding.
The source archive includes these materials and their build instructions.

\begin{thebibliography}{9}
\bibitem{oeis184} OEIS Foundation Inc., \emph{A281184}, entry by P.~D.~Hanna,
Jan.~17, 2017; conjecture by P.~Bala, Feb.~16, 2026; revision 41,
Nov.~6, 2026. \url{https://oeis.org/A281184}.
\bibitem{su} L.-H.~Su, \emph{Quadratic Derivations and Uniform Modular
Periods for Six Inverse-Function Sequences}, version 1, Oct.~31, 2026.
Theorem 3.2, Corollary 3.3, Propositions 3.4--3.6, Theorem 4.1 and
the following normalization remark.
\url{https://doi.org/10.5281/zenodo.23835577}.
\bibitem{graves} P.~Bala, \emph{The method of Graves for computing inverse
functions}, Jan.~26, 2026, Theorem 1 and Corollary 1.
\url{https://oeis.org/A281181/a281181.pdf}.
\bibitem{oeis428} OEIS Foundation Inc., \emph{A281428}, revision 54,
Nov.~6, 2026; P.~Bala's Jan.~22 and Mar.~3, 2026 blocks.
\url{https://oeis.org/A281428}.
\bibitem{oeis180} OEIS Foundation Inc., \emph{A281180}, revision 86,
Nov.~6, 2026. \url{https://oeis.org/A281180}.
\bibitem{oeis181} OEIS Foundation Inc., \emph{A281181}, revision 103,
Nov.~6, 2026. \url{https://oeis.org/A281181}.
\bibitem{oeis182} OEIS Foundation Inc., \emph{A281182}, revision 32,
Sept.~3, 2017. \url{https://oeis.org/A281182}.
\bibitem{oeis183} OEIS Foundation Inc., \emph{A281183}, revision 60,
Nov.~6, 2026. \url{https://oeis.org/A281183}.
\bibitem{oeis440} OEIS Foundation Inc., \emph{A281440}, revision 17,
Nov.~6, 2026. \url{https://oeis.org/A281440}.
\end{thebibliography}
\medskip\noindent\small
\copyright\ 2026 Haobo Ma and Wenlin Zhang.
Licensed under Creative Commons Attribution 4.0 International
(CC BY 4.0), \url{https://creativecommons.org/licenses/by/4.0/}.
\end{document}
